\sect{वैशाख-कृष्ण-वरूथिनी-एकादशी-माहात्म्यम्}
\label{sec:padma-vaishakha-krishna-varuthini ekadashi}

\ganapatyadiDhyanam
\padmaPuranam

\dnsub{अथ पञ्चाशत्तमोऽध्यायः॥५०}

\uvacha{युधिष्ठिर उवाच}

\twolineshloka
{वैशाखस्यासिते पक्षे किन्नामैकादशी भवेत्}
{महिमानं कथय मे वासुदेव नमोऽस्तु ते}% ॥१॥

\uvacha{श्रीकृष्ण उवाच}

\twolineshloka
{सौभाग्यदायिनी राजन्निहलोके परत्र च}
{वैशाखकृष्णपक्षे तु नाम्ना चैव वरूथिनी}% ॥२॥

\twolineshloka
{वरूथिन्या व्रतेनैव सौख्यं भवति सर्वदा}
{पापहानिश्च भवति सौभाग्यप्राप्तिरेव च}% ॥३॥

\twolineshloka
{दुर्भगा या करोत्येनां सा स्त्री सौभाग्यमाप्नुयात्}
{लोकानां चैव सर्वेषां भुक्तिमुक्तिप्रदायिनी}% ॥४॥

\twolineshloka
{सर्वपापहरा नॄणां गर्भवासनिकृन्तनी}
{वरूथिन्या व्रतेनैव मान्धाता स्वर्गतिं गतः}% ॥५॥

\twolineshloka
{धुन्धुमारादयश्चान्ये राजानो बहवस्तथा}
{ब्रह्मकपालनिर्मुक्तो बभूव भगवान् भवः}% ॥६॥

\threelineshloka
{दशवर्षसहस्राणि तपस्तप्यति यो नरः}
{कुरुक्षेत्रे  रविग्रहे स्वर्णभारं ददाति यः}
{तत्तुल्यं फलमाप्नोति वरूथिन्या व्रतं चरन्}% ॥७॥

\twolineshloka
{श्रद्धावान् यस्तु कुरुते वरूथिन्या व्रतं नरः}
{वाञ्छितं लभते सोऽपि इहलोके परत्र च}% ॥८॥

\twolineshloka
{पवित्रा पावनी ह्येषा महापातकनाशिनी}
{भुक्तिमुक्तिप्रदा चैव कर्तॄणां नृपसत्तम}% ॥९॥

\twolineshloka
{अश्वदानान्नृपश्रेष्ठ गजदानं विशिष्यते}
{गजदानाद् भूमिदानं तिलदानं ततोऽधिकम्}% ॥१०॥

\twolineshloka
{तस्माच्च स्वर्णदानं वै अन्नदानं ततोऽधिकम्}
{अन्नदानात् परं दानं न भूतं न भविष्यति}% ॥११॥

\twolineshloka
{पितृदेवमनुष्याणां तृप्तिरन्नेन जायते}
{तत्समं कविभिः प्रोक्तं कन्यादानं नृपोत्तम}% ॥१२॥

\twolineshloka
{धेनुदानं च तत्तुल्यमित्याह भगवान् स्वयम्}
{प्रोक्तेभ्यः सर्वदानेभ्यो विद्यादानं विशिष्यते}% ॥१३॥

\twolineshloka
{तत्फलं समवाप्नोति नरः कृत्वा वरूथिनीम्}
{कन्यावित्तेन जीवन्ति ये नराः पापमोहिताः}% ॥१४॥

\twolineshloka
{पुण्यक्षयं ते गच्छन्ति निरयं यातनामयम्}
{तस्मात् सर्वप्रयत्नेन न ग्राह्यं कन्यकाधनम्}% ॥१५॥

\twolineshloka
{यश्च गृह्णाति लोभेन कन्यां क्रीत्वा च तद् धनम्}
{सोऽन्यजन्मनि राजेन्द्र ओतुर्भवति निश्चितम्}% ॥१६॥

\twolineshloka
{कन्यां पुण्येन यो दद्याद् यथाशक्ति स्वलङ्कृताम्}
{तत् पुण्यसङ्ख्यां नृपते चित्रगुप्तो न शक्नुयात्}% ॥१७॥

\twolineshloka
{तत्तुल्यं फलमाप्नोति नरः कृत्वा वरूथिनीम्}
{कांस्यं मांसं मसूरांश्च चणकान् कोद्रवांस्तथा}% ॥१८॥

\twolineshloka
{शाकं मधु परान्नं च पुनर्भोजनमैथुने}
{वैष्णवो व्रतकर्ता च दशम्यां दश वर्जयेत्}% ॥१९॥

\twolineshloka
{द्यूतं क्रीडां च निद्रां च ताम्बूलं दन्तधावनम्}
{परापवादं पैशुन्यं स्तेयं हिंसां तथा रतिम्}% ॥२०॥

\twolineshloka
{क्रोधं चैवानृतं वाक्यमेकादश्यां विवर्जयेत्}
{कांस्यं मांसं सुरां क्षौद्रं  तैलं पतितभाषणम्}% ॥२१॥

\twolineshloka
{व्यायामं च प्रवासं च पुनर्भोजनमैथुनम्}
{वृषपृष्ठं मसूरान्नं द्वादश्यां परिवर्जयेत्}% ॥२२॥

\threelineshloka
{अनेन विधिना राजन् विहिता यैर्वरूथिनी}
% {सर्वपापक्षयं कृत्वा दद्यात् प्रान्तेऽक्षयां गतिम्}
{रात्रौ जागरणं कृत्वा पूजितो मधुसूदनः}% ॥२३॥

\twolineshloka
{सर्वपापविनिर्मुक्तास्ते यान्ति परमां गतिम्}
{तस्मात् सर्वप्रयत्नेन कर्तव्या पापभीरुभिः}% ॥२४॥

\threelineshloka
{क्षपारितनयाद् भीतो नरः कुर्याद् वरूथिनीम्}
{पठनाच्छ्रवणाद् राजन् गोसहस्रफलं लभेत्}
{सर्वपापविनिर्मुक्तो विष्णुलोके महीयते}% ॥२५॥

॥इति श्रीपाद्मे महापुराणे पञ्चपञ्चाशत्साहस्र्यां संहितायामुत्तरखण्डे उमापतिनारदसंवादान्तर्गतकृष्णयुधिष्ठिरसंवादे वैशाख-कृष्ण-वरूथिनी-एकादशी-माहात्म्यं नाम पञ्चाशत्तमोऽध्यायः॥५०॥

\genMangalaShloka

\hyperref[sec:ekadashi_mahatmyam_padma_puranam]{\closesub}
\clearpage

